\section{Learning LQR}

\subsection*{(a)}
The model-based recursive least-squares method uses the most structure. It
assumes the dynamics are linear, the cost is quadratic, and that once
\(A,B,Q,R\) are estimated, the controller can be found by solving the discounted
DARE. That is a lot to assume, but it also explains why it performs the best in
the plots: the gain error and cost drop quickly and stay close to optimal.

The model-free policy iteration method uses less information. It does not try to
identify \(A,B,Q,R\), but it still uses LQR structure by assuming the Q-function
is quadratic and the policy is linear. Because it uses less structure, it does
not perform as well as the model-based method, but it still does better than
direct policy gradient.

The model-free policy-gradient method uses the least problem information. It
only assumes a linear Gaussian policy and estimates the update from sampled
rollouts. This makes it more general, but the updates have higher variance, so
it is slower and less stable.

\subsection*{(b)}
Prior beliefs about \(A,B,Q,R\) fit most naturally into the model-based method.
They could be used as regularization terms, Bayesian priors, or just better
initial guesses for the least-squares estimates.

For model-free policy iteration, those priors are less direct because the method
learns the Q-function Hessian instead of \(A,B,Q,R\). Still, they could help
initialize the policy \(K\), choose regularization, or initialize the quadratic
Q-function estimate.

For policy gradient, the priors are the hardest to use directly. They could help
pick a better initial policy, tune the exploration covariance, or design a
baseline, but the algorithm itself never estimates \(A,B,Q,R\).
